% Part: incompleteness
% Chapter: representability-in-q
% Section: c-representable

\documentclass[../../../include/open-logic-section]{subfiles}

\begin{document}

\olfileid{inc}{req}{cre}
\olsection{$C$-এর অপেক্ষকগুলি $\Th{Q}$-তে উপস্থাপনযোগ্য}

দেখাতে হবে, $C$-এর প্রত্যেক অপেক্ষক $\Th{Q}$-তে উপস্থাপনযোগ্য। শেষ পর্যন্ত
$C$-এর প্রতিটি $k$-স্থানীয় অপেক্ষক $f(x_0,\dots,x_{k-1})$-কে এমন একটি
সূত্র $!A_f(x_0,\dots,x_{k-1},y)$ দেওয়ার উপায় দেখাতে হবে, যা তাকে
উপস্থাপন করে।

\begin{lem}
স্বাভাবিক সংখ্যা $n$ ও $m$ দেওয়া থাকলে, $n \neq m$ হলে $Q \Proves \num
n \neq \num m$।
\end{lem}

\begin{proof}
$n$-এর উপর আরোহ ব্যবহার করে দেখাই: প্রত্যেক $m$-এর জন্য $n \neq m$ হলে
$Q \Proves \eq/[\num n][\num m]$।

ভিত্তি ক্ষেত্রে $n = 0$। $m$, $0$-র সমান না হলে কোনো স্বাভাবিক সংখ্যা
$k$-এর জন্য $m = k + 1$। আমাদের একটি স্বতঃসিদ্ধ বলে
$\lforall[x][\eq/[0][x']]$। একটি পরিমাণসূচক-স্বতঃসিদ্ধে $x$-এর স্থানে
  $\num k$ বসিয়ে $\eq/[0][\num k']$ সিদ্ধান্তে আসতে পারি। কিন্তু
  $\num k'$-ই $\num m$।

আরোহের ধাপে দাবিটি $n$-এর জন্য সত্য ধরে $n+1$ বিবেচনা করি। $m$ যেকোনো
স্বাভাবিক সংখ্যা হোক। দুটি সম্ভাবনা: হয় $m = 0$, নয়তো কোনো $k$-এর জন্য
$m = k+1$। প্রথম ক্ষেত্রটি উপরের মতোই সামলানো যায়। দ্বিতীয় ক্ষেত্রে ধরি
$n+1 \neq k+1$। তাহলে $n \neq k$। $n$-এর জন্য আরোহের অনুমান থেকে পাই
$\Th{Q} \Proves \eq/[\num n][\num k]$। একটি স্বতঃসিদ্ধ বলে
$\lforall[x][\lforall[y][\eq[x'][y'] \lif \eq[x][y]]]$। একটি
পরিমাণসূচক-স্বতঃসিদ্ধ ব্যবহার করে পাই $\eq[\num n'][\num k'] \lif
\eq[\num n][\num k]$। বচনগত যুক্তিবিদ্যা ব্যবহার করে $\Th{Q}$-তে
$\eq/[\num n][\num k] \lif \eq/[\num n'][\num k']$ সিদ্ধান্তে আসতে
পারি। মোডাস পোনেন্স থেকে পাই $\eq/[\num n'][\num k']$; এটিই প্রয়োজন,
কারণ $\num k'$ হলো $\num m$।
\end{proof}

লেমাটি খুব বেশি কিছু বলে না: মূলত $\Th{Q}$ যে আলাদা সংখ্যাপদ আলাদা
বস্তুকে নির্দেশ করে—এটি প্রমাণ করতে পারে, সেটিই বলে। যেমন, $\Th{Q}$
$0'' \neq 0'''$ প্রমাণ করে। কিন্তু সাধারণভাবে কথাটি দেখাতে কিছু
সতর্কতা দরকার। আরও লক্ষ করো, আমরা আরোহ ব্যবহার করলেও সেটি
$\Th{Q}$-এর \emph{বাইরের} আরোহ।

শূন্য, উত্তরসূরি, যোগ, গুণ, সমতার চরিত্রসূচক অপেক্ষক এবং অভিক্ষেপগুলিকে
আমরা উপস্থাপন করতে পারব। প্রতিটি ক্ষেত্রেই উপযুক্ত উপস্থাপক সূত্র
একেবারেই সরল; যেমন, শূন্যকে উপস্থাপন করে
\[
y = 0,
\]
উত্তরসূরিকে উপস্থাপন করে
\[
x_0' = y,
\]
এবং যোগকে উপস্থাপন করে
\[
x_0 + x_1 = y.
\]
আসল কাজ হলো দেখানো যে $\Th{Q}$ প্রাসঙ্গিক বক্তব্যগুলি প্রমাণ করতে পারে।
যেমন, উপরের সূত্র যোগকে উপস্থাপন করে বলার জন্য দেখাতে হয় যে প্রত্যেক জোড়া
স্বাভাবিক সংখ্যা $m$ ও $n$-এর জন্য $\Th{Q}$ প্রমাণ করে
\[
\eq[\num n + \num m][\num {n+m}]
\]
এবং
\[
\lforall[y][(\eq[(\num n + \num m)][y] \lif \eq[y][\num {n+m}])].
\]

মিশ্রণের ক্ষেত্রে কী হবে? ধরি $h$-কে এভাবে সংজ্ঞায়িত করা হয়েছে:
\[
h(x_0,\dots,x_{l-1}) = f(g_0(x_0,\dots,x_{l-1}), \dots,
g_{k-1}(x_0,\dots,x_{l-1})).
\]
এখানে $f,g_0,\dots,g_{k-1}$ অপেক্ষকগুলিকে যথাক্রমে উপস্থাপন করে
এমন সূত্র $!A_f, !A_{g_0}, \dots,
!A_{g_{k-1}}$ আমরা ইতিমধ্যে পেয়েছি। তাহলে $h$-কে উপস্থাপন করে এমন
একটি সূত্র~$!A_h$ সংজ্ঞায়িত করা যায়: $!A_h(x_0,\dots,x_{l-1},y)$
হবে
\begin{multline*}
  \lexists[z_0][\dots \lexists[z_{k-1}][(!A_{g_0}(x_0,\dots,x_{l-1},z_0) \land
  \dots \land {}\\
  !A_{g_{k-1}}(x_0,\dots,x_{l-1},z_{k-1}) \land !A_f(z_0,\dots,z_{k-1},y))]]
\end{multline*}
% BN-SRC-503 TARGET-CORRECTION-BEGIN
\par\smallskip\noindent\textnormal{\small\textbf{উৎস-সংশোধন (BN-SRC-503):} মূল মিশ্রণ-সূত্রে প্রথম অস্তিত্বসূচকের চলক-ঘরে \texttt{z\_0\textbackslash dots} বসেছে, মাঝখানে বন্ধনী অকালে শেষ হয়েছে এবং শেষের উপস্থাপক সূত্রটি পরিমাণসূচকের পরিসরের বাইরে পড়েছে। বাংলা সূত্রে $z_0$ থেকে $z_{k-1}$ পর্যন্ত সঠিক পরিসর ও বন্ধনী রাখা হয়েছে; বর্তমান অধ্যায়ের মিশ্রণ-বিভাগেও সমতুল্য সংশোধন আছে।}
% BN-SRC-503 TARGET-CORRECTION-END

শেষে সীমাহীন অনুসন্ধান বিবেচনা করি। ধরি $g(x,\vec z)$ নিয়মিত এবং
$\Th{Q}$-তে উপস্থাপনযোগ্য; ধরা যাক উপস্থাপক সূত্রটি $!A_g(x,\vec
z,y)$। $f(\vec z) = \mu x \; g(x,\vec z)$ দিয়ে $f$ সংজ্ঞায়িত হোক।
$f$-কে উপস্থাপনকারী একটি সূত্র $!A_f(\vec z,y)$ খুঁজতে চাই। স্বাভাবিক
পছন্দটি হলো:
\[
!A_f(\vec z,y) \ident !A_g(y,\vec z,0) \land \lforall[w][(w < y
\lif \lnot !A_g(w,\vec z,0))].
\]

আরও কয়েকটি লেমার সাহায্যে দেখানো যায় যে এই সূত্রটি কাজ করে। যেমন,

\begin{lem}
  প্রত্যেক চলক $x$ ও স্বাভাবিক সংখ্যা $n$-এর জন্য $\Th{Q}$ প্রমাণ করে
  $\eq[(x' + \num n)][(x + \num n)']$।
\end{lem}

এটি যে $\Th{Q}$-তে $\lforall[x][\lforall[y][(x' + y) = (x +y)']]$
প্রমাণযোগ্য বলা অপেক্ষা দুর্বল (শেষোক্তটি মিথ্যা), তা আবারও লক্ষ করো।

\begin{proof}
যথারীতি $n$-এর উপর আরোহ দিয়ে প্রমাণ করি। ভিত্তি ক্ষেত্রে $n =
0$, দেখাতে হবে $\Th{Q}$ প্রমাণ করে $\eq[(x'+0)][(x+0)']$। কিন্তু
আমাদের আছে:
\begin{eqnarray*}
& & \eq[(x' + 0)][x'] \quad \text{স্বতঃসিদ্ধ ৪ থেকে} \\
& & \eq[(x + 0)][x] \quad \text{স্বতঃসিদ্ধ ৪ থেকে} \\
& & \eq[(x+0)'][x'] \quad \text{দ্বিতীয় পংক্তি থেকে} \\
& & \eq[(x' + 0)][(x + 0)'] \quad \text{প্রথম ও তৃতীয় পংক্তি থেকে}
\end{eqnarray*}
আরোহের ধাপে $\Th{Q}$-তে $\eq[(x' +
  \num n)][(x + \num n)']$ নিষ্পাদিত হয়েছে বলে ধরি। যেহেতু
$\num{n+1}$ হলো $\num n'$, দেখাতে হবে $\Th{Q}$ প্রমাণ করে
$\eq[(x' + \num n')][(x + \num n')']$। নিচের সমতাগুলি
$\Th{Q}$-তে নিষ্পাদন করা যায়:
\begin{eqnarray*}
(x' + \num n') & \eq & (x' + \num n)' \quad \text{স্বতঃসিদ্ধ ৫} \\
& \eq & (x + \num n')' \quad \text{আরোহের অনুমান থেকে}
\end{eqnarray*}
\end{proof}

\begin{lem}
\begin{enumerate}
\item $\Th{Q}$ প্রমাণ করে $\lnot (x < \num 0)$।
\item প্রত্যেক স্বাভাবিক সংখ্যা $n$-এর জন্য $\Th{Q}$ প্রমাণ করে
\[
x < \num {n+1} \lif (\eq[x][\num 0] \lor \dots \lor \eq[x][\num n]).
\]
\end{enumerate}
\end{lem}

\begin{proof}
প্রথম দফাটি এবং দ্বিতীয়টির একাংশ অনানুষ্ঠানিকভাবে করি—অর্থাৎ আনুষ্ঠানিক
!!{derivation} নির্মাণের জন্য শুধু ইঙ্গিত দিই।

প্রথম দফায়, $<$-এর সংজ্ঞা অনুসারে $\Th{Q}$-তে $\lnot
\lexists[y][\eq[(y'+x)][0]]$ প্রমাণ করতে হবে। প্রথম-ক্রম যুক্তিবিদ্যার
স্বতঃসিদ্ধ ও বিধি অনুসারে এটি $\lforall[y][\eq/[(y' +
    x)][0]]$-এর সমতুল্য। ধারণাটি এই: ধরি $\eq[(y'+x)][0]$। $x$ শূন্য
হলে $\eq[(y' + 0)][0]$ পাই। কিন্তু $\Th{Q}$-র চতুর্থ স্বতঃসিদ্ধে
$\eq[(y' + 0)][y']$; আর দ্বিতীয় স্বতঃসিদ্ধে $\eq/[y'][0]$—স্ববিরোধ।
তাই $\lforall[y][\eq/[(y' +x)][0]]$। $x$, $0$ না হলে তৃতীয়
স্বতঃসিদ্ধ অনুসারে এমন একটি $z$ আছে যেন $\eq[x][z']$। তখন
$\eq[(y' + z')][0]$ পাই। পঞ্চম স্বতঃসিদ্ধে $\eq[(y' + z)'][0]$ পাই;
এটিও দ্বিতীয় স্বতঃসিদ্ধের বিরোধী।

দ্বিতীয় দফায় $n$-এর ওপর আরোহ করি। ভিত্তি ক্ষেত্রে $n =0$ বিবেচনা
করি। তখন $x < \num 1 \lif \eq[x][\num
  0]$ দেখাতে হবে। ধরি $x < \num 1$। তাহলে $<$-এর সংজ্ঞাদায়ী
স্বতঃসিদ্ধ অনুসারে $\lexists[y][\eq[(y'+x)][0']]$। $y$-এর ওই ধর্ম
আছে ধরি; তখন $\eq[y'+x][0']$।

দেখাতে হবে $\eq[x][0]$। তৃতীয় স্বতঃসিদ্ধ অনুসারে, $x$ শূন্য না হলে
কোনো $z$-এর জন্য তা $z'$-এর সমান। তখন $\eq[(y' + z')][0']$। $\Th{Q}$-র পঞ্চম
স্বতঃসিদ্ধে $\eq[(y'+z)'][0']$। প্রথম স্বতঃসিদ্ধে $\eq[(y' +
    z)][0]$। কিন্তু সংজ্ঞা অনুসারে এর অর্থ $z < 0$, যা প্রথম দফার
বিরোধী।
\end{proof}

\begin{explain}
আমরা দেখিয়েছি, গণনসাধ্য অপেক্ষকের সেটকে $\Th{Q}$-তে উপস্থাপনযোগ্য
অপেক্ষকের সেট হিসেবে চিহ্নিত করা যায়। বস্তুত প্রমাণটি আরও সাধারণ।
উপস্থাপনযোগ্যতার সংজ্ঞা থেকে সহজেই দেখা যায়, $\Th{Q}$-কে প্রসারিত করে
এমন (অথবা যার মধ্যে $\Th{Q}$-কে ব্যাখ্যা করা যায় এমন) যেকোনো সঙ্গত
তত্ত্ব গণনসাধ্য অপেক্ষকগুলিকে উপস্থাপন করতে পারে। বিপরীত দিকে, যে
সঙ্গতিপূর্ণ !!{derivation} পদ্ধতিতে প্রমাণের ধারণাটি গণনসাধ্য, সেখানে
প্রত্যেক উপস্থাপনযোগ্য অপেক্ষক গণনসাধ্য। তাই, যেমন, গণনসাধ্য অপেক্ষকের
সেটকে পেয়ানো পাটিগণিত, এমনকি জার্মেলো--ফ্রেঙ্কেল সেটতত্ত্বে
উপস্থাপনযোগ্য অপেক্ষকের সেট হিসেবেও চিহ্নিত করা যায়। গ্যোডেল লক্ষ
করেছিলেন, এটি কিছুটা বিস্ময়কর। প্রমাণযোগ্যতার ক্ষেত্রে কোন তত্ত্ব
বিবেচনা করছি, প্রশ্নগুলি তার প্রতি খুব সংবেদনশীল—এ কথা আমরা দেখব;
মোটামুটিভাবে, স্বতঃসিদ্ধ যত শক্তিশালী, তত বেশি প্রমাণ করা যায়। কিন্তু
সঙ্গতিপূর্ণ স্বতঃসিদ্ধায়িত তত্ত্বের এক বিস্তৃত শ্রেণি জুড়ে উপস্থাপনযোগ্য
অপেক্ষক ঠিক গণনসাধ্য অপেক্ষকগুলিই; শক্তিশালী তত্ত্ব বেশি অপেক্ষক
উপস্থাপন করে না।
% BN-SRC-504 TARGET-CORRECTION-BEGIN
\par\smallskip\noindent\textnormal{\small\textbf{উৎস-সংশোধন (BN-SRC-504):} মূল সাধারণীকরণে কেবল গণনযোগ্য প্রমাণপদ্ধতির কথা আছে; অসঙ্গত তত্ত্বে যেকোনো সূত্র প্রমাণযোগ্য বলে উপস্থাপনযোগ্যতার সংজ্ঞা অগণনসাধ্য অপেক্ষকের জন্যও তুচ্ছভাবে পূর্ণ হয়। বাংলা পাঠে সঙ্গতির শর্ত যোগ করা হয়েছে, যেমন বর্তমান অধ্যায়ের অনুরূপ অনুচ্ছেদেও করা হয়েছে।}
% BN-SRC-504 TARGET-CORRECTION-END
\end{explain}

\end{document}
